% from ega4-16-fr.tex
\oldpage[IV-4]{12}
déduit de l'homomorphisme $\nu_n$ (16.2.1.2) est bijectif. En outre,
l'homomorphisme de $\mathcal{O}_{Y'}$-Modules
\[
  \operatorname{Gr}_1(u):
  \mathscr{G}r_1(f)\otimes_{\mathcal{O}_Y}\mathcal{O}_{Y'}
  \longrightarrow\mathscr{G}r_1(f')
  \tag{16.2.3.2}
  \label{IV.16.2.3.2}
\]
est bijectif.
\end{enumerate}
\end{corollary}

\noindent(i) Notons d'abord que les morphismes $f':Y'\to X'$ et $u:Y'\to Y$
identifient $Y'$ au produit $Y\times_X X'$ (pour les morphismes structuraux
$f:Y\to X$ et $v:X'\to X$) (14.5.12.1). La conclusion de (i) résulte alors de
(16.2.2, (ii)), le morphisme $g$ étant une immersion.

\noindent(ii) Le diagramme commutatif
\[
  \xymatrix{
    Y_f^{(n)} \ar[d]^{h_n}
      & Y_{f'}'^{(n)} \ar[d]^-{h'_n} \ar[l]^{w_n} \\
    X \ar[d]^g
      & X' \ar[d]^{g'} \ar[l]^v \\
    Y
      & Y' \ar[l]^u
  }
\]
identifie $Y_{f'}'^{(n)}$ au produit $Y_f^{(n)}\times_X X'$ et $X'$ au produit
$X\times_Y Y'$, donc (I, 3.3.9) il identifie (pour les morphismes
$g'\circ h'_n$ et $w_n$) $Y_{f'}'^{(n)}$ au produit
$Y_f^{(n)}\times_Y Y'$. Comme $Y_f^{(n)}$ (resp. $Y_{f'}'^{(n)}$) est le
préschéma affine au-dessus de $Y$ (resp. $Y'$) associé à la
$\mathcal{O}_Y$-Algèbre $\mathcal{O}_{Y_f^{(n)}}$ (resp. à la
$\mathcal{O}_{Y'}$-Algèbre $\mathcal{O}_{Y_{f'}'^{(n)}}$), le fait que
l'homomorphisme canonique (16.2.3.1) soit bijectif résulte de (II, 1.5.2).
Enfin, l'homomorphisme canonique (16.2.3.1) est compatible avec les
augmentations
$\mathcal{O}_{Y_f^{(n)}}\to\mathcal{O}_Y$ et
$\mathcal{O}_{Y_{f'}'^{(n)}}\to\mathcal{O}_{Y'}$; comme
$\mathcal{O}_{Y_f^{(n)}}$ est somme directe (en tant que
$\mathcal{O}_Y$-Module) de $\mathcal{O}_Y$ et de l'Idéal d'augmentation
$\mathscr{I}/\mathscr{I}^{n+1}$, on voit donc que l'homomorphisme canonique
(16.2.3.1), restreint à
$(\mathscr{I}/\mathscr{I}^{n+1})\otimes_{\mathcal{O}_Y}\mathcal{O}_{Y'}$,
est une bijection de ce dernier sur
$\mathscr{I}'/\mathscr{I}'^{n+1}$. Pour $n=1$, cela montre que
$\operatorname{Gr}_1(u)$ est bijectif.

On notera que, sous les hypothèses de (16.2.3), les homomorphismes
$\operatorname{Gr}_n(u)$ sont surjectifs en vertu de ce qui précède, mais ne
sont pas bijectifs en général pour $n\geq 2$. Toutefois:

\begin{corollary}[16.2.4]
\label{IV.16.2.4}
Sous les hypothèses de (16.2.3), supposons que $u:Y'\to Y$ soit un morphisme
plat (resp. que les $\mathscr{G}r_n(f)$ soient des $\mathcal{O}_Y$-Modules
plats pour $n\leq m$). Alors l'homomorphisme
\[
  \operatorname{Gr}_n(u):
  \mathscr{G}r_n(f)\otimes_{\mathcal{O}_Y}\mathcal{O}_{Y'}
  \longrightarrow\mathscr{G}r_n(f')
\]
est bijectif pour tout $n$ (resp. pour $n\leq m$).
\end{corollary}

Si $u$ est plat, il en est de même de $v:X'\to X$ qui s'en déduit par
changement de base, et on sait déjà dans ce cas que $\operatorname{Gr}(u)$ est
bijectif (16.2.2, (iii)). Si les $\mathscr{G}r_n(f)$ sont plats pour
$n\leq m$, on voit d'abord par récurrence sur $n$ qu'il en est de même des
$\mathscr{I}/\mathscr{I}^{n+1}$ pour $n\leq m$, en vertu des suites exactes
\[
  0\longrightarrow\mathscr{I}^n/\mathscr{I}^{n+1}
  \longrightarrow\mathscr{I}/\mathscr{I}^{n+1}
  \longrightarrow\mathscr{I}/\mathscr{I}^n
  \longrightarrow 0
\]
